\inicializarSeccion{Marco Teórico}

\inicializarSubseccion{Frenado magnetico}

\tab El frenado magnético en una torre (góndola) de caída libre está relacionado con distintos principios del electromagnetismo. Los conceptos más importantes serán las leyes de Faraday, Lenz y Biot-Savart además de las corrientes de Eddy; estas ayudan a describir cómo es que se comporta un campo magnético junto a las fuerzas que se crean durante el proceso de frenado.

El frenado magnético ayuda a crear un sistema de frenado para disminuir la velocidad de un objeto usando campos magnéticos que permiten que el frenado se realice sin contacto físico. Gracias a la implementación de este método de frenado, se puede reducir el desgaste de los materiales y se puede crear un frenado más suave con la góndola.

\inicializarSubseccion{Ley de las corrientes de Eddy}

\tab Las corrientes de Eddy son corrientes eléctricas inducidas en los materiales conductores cuando atraviesan un campo magnético variable. En nuestro caso, cuando la góndola cae genera corrientes inducidas, las cuales generan un campo magnético que produce una fuerza de frenado con la cual la velocidad irá disminuyendo gradualmente.

\inicializarSubseccion{Ley de Faraday}

\tab La ley de Faraday propone un cambio en el flujo magnético producido por una fuerza electromotriz. Con esta ley podemos entender cómo es que se crean las corrientes inducidas en una góndola en movimiento dentro de un campo magnético.
La fórmula que se utiliza es la siguiente:

\begin{equation}
    \varepsilon = -\frac{d\Phi_B}{dt}
\end{equation}

Donde $\varepsilon$ es la fuerza inducida y $\Phi_B$ es el flujo magnético.

\inicializarSubseccion{Ley de Lenz}

\tab La ley de Lenz nos habla de que un campo magnético generado por una corriente inducida se opondrá a la fuerza que la creó. Gracias a ello se genera la fuerza que desacelera la góndola en la torre.

\inicializarSubseccion{Ley de Biot-Savart}

\tab La ley de Biot-Savart nos ayuda a calcular el campo magnético que se genera por medio de una corriente eléctrica. Para poder hacer el cálculo utilizamos la siguiente fórmula:

\begin{equation}
    \vec{B} = \frac{\mu_0}{4\pi} \int \frac{I d\vec{l} \times \hat{r}}{r^2}
\end{equation}

$I$ representa la corriente eléctrica, $r$ es la distancia, $\mu_0$ es la permeabilidad y $d\vec{l}$ es el diferencial. Esta ley facilita el modelado en MATLAB de nuestro proyecto.

\inicializarSubseccion{Campo magnetico}

\tab El campo magnético en una línea de corriente ocurre cuando circula una corriente eléctrica por un conductor recto, eso crea un campo alrededor del conductor junto a una magnitud del mismo campo. Esto lo podemos calcular por medio de la siguiente fórmula:

\begin{equation}
    B = \frac{\mu_0 I}{2\pi r}
\end{equation}

Donde $B$ es el campo magnético, $I$ es la corriente eléctrica y $r$ es la distancia que existe hacia el conductor.

Con el resultado que se obtiene podemos observar que el campo disminuye conforme la distancia al conductor se hace más grande.

\inicializarSubseccion{Espira circular}

\tab También tenemos el campo magnético en una espira circular, el cual es un conductor que tiene forma de una circunferencia la cual tiene una corriente eléctrica. Debido a sus características, se genera un campo magnético que se concentra en el centro de esta espira.
Para poder calcular este campo magnético utilizamos la siguiente fórmula:

\begin{equation}
    B = \frac{\mu_0 I}{2R}
\end{equation}

Donde $B$ es el campo magnético, $I$ la corriente eléctrica, $R$ es el radio que tiene la espira.

\inicializarSubseccion{Simulación de MATLAB}

\tab Gracias a estos conceptos podemos entender cómo es que funciona el frenado magnético de una góndola, con esto somos capaces de generar un modelo por medio de MATLAB donde mostramos el funcionamiento de estas leyes aplicadas en nuestra góndola.

\inicializarSubseccion{Momento de dipolo magnético}

\tab El momento de dipolo magnético es un modelo que representa el comportamiento de imán, el cual se describe el momento dipolar magnético, el cual indica la intensidad y dirección del imán.

\begin{equation}
    \vec{\mu} = -m\hat{\mathbf{k}}
\end{equation}

\inicializarSubseccion{Campo magnético sobre el eje de una espira circular}

\tab El campo magnético es generado por una espira circular con una corriente puede calcularse sobre el eje de la z, en este caso el campo apunta en la dirección del eje z.

\begin{equation}
    B_z = \frac{\mu_0 I R^2}{2(R^2 + z^2)^{3/2}}
\end{equation}

Donde $\mu_0$ es la permeabilidad magnética, $I$ es la corriente eléctrica, $R$ es el radio de la espira y $z$ es la posición del eje de la espira.

\inicializarSubseccion{Fuerzas sobre un dipolo magnético}

\tab La fuerza que experimenta un dipolo magnético dentro de un campo magnético no uniforme se calcula con:

\begin{equation}
    \vec{F} = \nabla (\vec{\mu} \cdot \vec{B})
\end{equation}

Esta ecuación nos dice que el imán experimenta una fuerza cuando se encuentra en una región donde el campo magnético varía con la posición.

Sustituyendo el momento dipolar y el campo magnético de la espira obtenemos 

\begin{equation}
    \vec{F} = \frac{6\mu_0 I m R^2}{2(R^2 + z^2)^{5/2}} \hat{\mathbf{z}}
\end{equation}

Esta nos ayuda en la simulación de MATLAB, ya que la fuerza también puede aproximarse mediante la deriva del campo magnético respecto de la posición:

\begin{equation}
    \vec{F} = m\frac{dB_z}{dz} \hat{\mathbf{k}}
\end{equation}

Esta fuerza se suma vectorialmente con la fuerza de gravedad para encontrar la fuerza total que actúa sobre el imán durante el momento en el cual experimenta su caída.

\inicializarSubseccion{Flujo magnético a través de una espira}

\tab El flujo magnético representa la cantidad de campo magnético que atraviesa el área de una superficie. En el proyecto, el flujo del imán cambia conforme el imán se acerca o se aleja de la espira.

\begin{equation}
    \Phi_B = \int \vec{B} \cdot d\vec{A}
\end{equation}

Donde $\Phi_B$ es el flujo magnético, $\vec{B}$ es el campo magnético y $d\vec{A}$ es el diferencial del área.


\inicializarSubseccion{Fuerza electromotriz inducida}

\tab La ley de Faraday nos dice que una variación del flujo magnético genera una fuerza electromotriz inducida (FEM).

\begin{equation}
    \varepsilon = -\frac{d\Phi_B}{dt}
\end{equation}

Donde $\varepsilon$ es la fuerza electromotriz inducida.

\inicializarSubseccion{Movimiento rectilíneo no uniforme (MRUA)}

\tab Nos ayuda a calcular la posición del dipolo magnético durante la caída para conocer cómo es que va cambiando el flujo magnético con el tiempo.

\begin{equation}
    z(t) = z_0 + v_0 t + \frac{1}{2} a t^2
\end{equation}

Donde $z(t)$ es la posición del imán en función del tiempo, $z_0$ es la posición inicial, $v_0$ es la velocidad inicial y $a$ es la aceleración.

\inicializarSubseccion{Regla de la cadena aplicada a una FEM}

\tab Como el flujo magnético depende de la posición del imán y dicha posición va cambiando con el tiempo, se utiliza la regla de la cadena para calcular la FEM inducida.

\begin{equation}
    \varepsilon = -\frac{d\Phi_B}{dt} = -\frac{d\Phi_B}{dz} \cdot \frac{dz}{dt}
\end{equation}

Donde $\frac{dz}{dt}$ es la velocidad del imán.

\inicializarSubseccion{Método de Romberg para integración numérica}

\tab Este método de integración numérica nos ayuda a aproximar integrales definidas con mucha precisión a partir de mejoras sucesivas de la regla del trapecio. Optamos por utilizar ya que nos permite hacer mediciones con muchisimo menos error, sin tener que computar muchisimos puntos de la función a integrar, lo cual es de extrema importancia para poder obtener resultados precisos en nuestra simulación.

\begin{equation}
    I_{j,k}\approx \frac{4^k I_{j+1,k-1} - I_{j,k-1}}{4^k - 1}
\end{equation}


\textbf{NOTA}: Para la realización de este marco teórico se consultaron distintas fuentes, como lo son los videos y materiales proporcionados en el portal de Canvas, así como también libros y artículos relacionados con el tema. Entre las fuentes más importantes se encuentran \cite{griffiths2013introduction} y \cite{hayt2006teoría}, los cuales nos ayudaron a entender mejor el funcionamiento del frenado magnético y los conceptos físicos involucrados en el reto. Además, se consultaron otras fuentes para complementar la información y obtener una visión más completa del tema.